\documentclass[a4paper,11pt]{article}
\usepackage{cmap}
\usepackage[T2A]{fontenc}
\usepackage[utf8]{inputenc}
\usepackage[russian]{babel}
\usepackage{amsmath,amssymb,mathtools}
\usepackage{geometry}
\geometry{left=22mm,right=22mm,top=22mm,bottom=21mm,headheight=22pt,headsep=6mm,footskip=12mm}
\usepackage{microtype}
\usepackage{xcolor}
\usepackage{tikz}
\usetikzlibrary{arrows.meta,positioning,shapes.geometric,calc}
\usepackage{booktabs,tabularx,array}
\usepackage{enumitem}
\usepackage{fancyhdr}
\usepackage{setspace}
\usepackage{hyperref}
\usepackage{parskip}
\definecolor{accent}{HTML}{007C79}
\definecolor{darkaccent}{HTML}{005A58}
\definecolor{bodytext}{HTML}{202927}
\definecolor{muted}{HTML}{606E6A}
\definecolor{softgray}{HTML}{D0D9D6}
\color{bodytext}
\hypersetup{colorlinks=true,urlcolor=accent,linkcolor=accent,pdfauthor={Лёвин Артём Александрович},pdftitle={Чек-лист: координатный метод в стереометрии — Матвей Горбачев}}
\setstretch{1.16}
\setlength{\parindent}{0pt}
\setlist[enumerate]{leftmargin=8mm,itemsep=3mm,topsep=2mm}
\setlist[itemize]{leftmargin=6mm,itemsep=1.6mm,topsep=2mm}
\newcommand{\checksq}{\ensuremath{\square}\,}
\newcommand{\heading}[1]{\par\vspace{4mm}{\Large\color{darkaccent}#1}\par\vspace{1.7mm}{\color{accent}\rule{\linewidth}{0.55pt}}\par\vspace{2mm}}
\newcommand{\subheading}[1]{\par\vspace{2.2mm}{\bfseries\color{darkaccent}#1}\par\vspace{1.0mm}}
\newcommand{\key}[1]{{\color{darkaccent}\bfseries #1}}
\newcommand{\headername}{Координатный метод в стереометрии}
\pagestyle{fancy}
\fancyhf{}
\fancyhead[L]{\footnotesize\color{muted}\headername}
\fancyhead[R]{\footnotesize\color{muted}Матвей Горбачев}
\fancyfoot[C]{\footnotesize\color{muted}\thepage}
\renewcommand{\headrulewidth}{0.45pt}
\renewcommand{\headrule}{\hbox to\headwidth{\color{accent}\leaders\hrule height \headrulewidth\hfill}}
\newcommand{\cubesketch}[1]{%
\begin{tikzpicture}[scale=.82,line cap=round,line join=round,>=Latex]
\coordinate (A) at (0,0);\coordinate (B) at (3.1,0);\coordinate (D) at (1.0,1.05);\coordinate (C) at (4.1,1.05);
\coordinate (Aone) at (0,2.7);\coordinate (Bone) at (3.1,2.7);\coordinate (Done) at (1.0,3.75);\coordinate (Cone) at (4.1,3.75);
\draw[softgray,densely dashed] (A)--(D)--(Done) (D)--(C);
\draw[gray!75,line width=.75pt] (A)--(B)--(C)--(Cone)--(Done)--(Aone)--(A)--(B) (Aone)--(Bone)--(Cone) (Bone)--(B);
\draw[gray!75,line width=.75pt] (Aone)--(Done);
\draw[->,accent,line width=.65pt] (A)--(-.2,3.5) node[above,font=\scriptsize] {$z$};
\draw[->,accent,line width=.8pt] (A)--(3.8,0) node[right,font=\scriptsize] {$x$};
\draw[->,accent,line width=.8pt] (A)--(1.42,1.49) node[above,font=\scriptsize] {$y$};
\node[below left,font=\scriptsize] at (A) {$A$};
\node[below,font=\scriptsize] at (B) {$B$};
\node[left,font=\scriptsize] at (Aone) {$A_1$};
\node[above,font=\scriptsize] at (Cone) {$C_1$};
#1
\end{tikzpicture}}
\begin{document}
\thispagestyle{empty}
\vspace*{20mm}
\begin{center}
{\large\color{darkaccent} МАТЕМАТИКА}\par
\vspace{18mm}
{\fontsize{24.88}{30}\selectfont Чек-лист}\par
\vspace{12mm}
{\LARGE\color{accent} Координатный метод\\[2mm] в стереометрии}\par
\vspace{11mm}
{\large Уравнение плоскости и угол между плоскостями}\par
\vspace{25mm}
{\large Матвей Горбачев}\par
\vspace{9mm}
{\large 19 ноября 2026 г.}\par
\vfill
{Лёвин Артём Александрович}\par
\vspace{1mm}
{эксклюзивно для Матвея Горбачёва}\par
\vspace{18mm}
\begin{tikzpicture}
\draw[accent,line width=.8pt] (0,0) .. controls (1.1,.45) and (2.1,-.45) .. (3.3,0) .. controls (4.4,.45) and (5.5,-.45) .. (6.7,0) .. controls (7.8,.45) and (8.7,-.38) .. (10,0);
\fill[darkaccent] (3.3,0) circle (1.5pt);
\end{tikzpicture}
\end{center}
\clearpage

\heading{1. Как задать координаты в кубе}
\key{Главная идея.} Выберите начало координат в вершине $A$. Направьте оси вдоль трёх попарно перпендикулярных рёбер: $AB$, $AD$ и $AA_1$.

\begin{center}
\cubesketch{%
\draw[accent,line width=1.5pt] (A)--(B);
\draw[accent,line width=1.5pt] (A)--(Aone);
\draw[accent,line width=1.5pt] (A)--(D);
\node[below,font=\scriptsize] at (2.0,0) {$AB=a$};
}
\end{center}
При ребре $a$ получаем:
\[
A(0;0;0),\quad B(a;0;0),\quad D(0;a;0),\quad A_1(0;0;a),\quad C_1(a;a;a).
\]
\subheading{Точки на рёбрах}
Середина отрезка с концами $P(x_1;y_1;z_1)$ и $Q(x_2;y_2;z_2)$ имеет координаты
\[
K\!\left(\frac{x_1+x_2}{2};\frac{y_1+y_2}{2};\frac{z_1+z_2}{2}\right).
\]
Например, середина $BB_1$ имеет координаты $M\!\left(a;0;\frac a2\right)$.

\heading{2. Что означает уравнение плоскости}
Общее уравнение плоскости в пространстве:
\[
\boxed{Ax+By+Cz+D=0},\qquad A^2+B^2+C^2>0.
\]
Числа $A$, $B$, $C$ образуют \key{нормальный вектор} $\vec n=(A;B;C)$: он перпендикулярен плоскости. При умножении всех четырёх коэффициентов на одно ненулевое число плоскость сохраняется.

\begin{itemize}
\item Если плоскость проходит через начало координат $O(0;0;0)$, то $D=0$.
\item Если точка принадлежит плоскости, её координаты удовлетворяют уравнению.
\item Три точки определяют единственную плоскость, если они не лежат на одной прямой.
\end{itemize}

\clearpage
\heading{3. Как найти плоскость через три точки}
\key{Маршрут:} составьте $Ax+By+Cz+D=0$ $\longrightarrow$ подставьте координаты известных точек $\longrightarrow$ решите систему $\longrightarrow$ проверьте все три точки.

\subheading{Пример 1. Плоскость $AMC_1$}
В единичном кубе рассмотрим точки
\[
A(0;0;0),\qquad M\left(1;0;\frac12\right),\qquad C_1(1;1;1),
\]
где $M$ --- середина ребра $BB_1$.

\begin{center}
\cubesketch{%
\coordinate (M) at (3.1,1.35);
\draw[accent!80!black,line width=1.45pt] (A)--(M)--(Cone)--cycle;
\fill[accent] (M) circle (1.7pt) node[right,font=\scriptsize] {$M$};
}
\end{center}
Ищем уравнение $Ax+By+Cz+D=0$. Подстановка координат даёт
\[
\begin{cases}
D=0 & \text{(точка $A$)},\\[1mm]
A+\frac C2+D=0 & \text{(точка $M$)},\\[1mm]
A+B+C+D=0 & \text{(точка $C_1$)}.
\end{cases}
\]
Из первого равенства $D=0$, а из второго $A=-\frac C2$. Тогда из третьего
\[
-\frac C2+B+C=0\quad\Longrightarrow\quad B=-\frac C2.
\]
Берём $C=-2$ (коэффициенты можно умножить на общий множитель). Получаем $A=B=1$.
\[
\boxed{\alpha:\ x+y-2z=0},\qquad \boxed{\vec n_\alpha=(1;1;-2)}.
\]
\subheading{Проверка обязательна}
\[
A:\ 0=0,\quad M:\ 1+0-2\cdot\tfrac12=0,\quad C_1:\ 1+1-2=0.
\]
\key{Обратите внимание:} у $y$ стоит знак «плюс». Смена этого знака нарушила бы условие принадлежности $C_1$ плоскости.

\clearpage
\heading{4. Вторая плоскость: свободный член $D$}
\subheading{Пример 2. Плоскость $MTQ$}
Пусть координаты трёх точек известны:
\[
M(2;0;1),\qquad T(1;0;0),\qquad Q(0;1;0).
\]
В кубе с ребром $2$ это середины рёбер $BB_1$, $AB$ и $AD$ соответственно.
\begin{center}
\cubesketch{%
\coordinate (M) at (3.1,1.35);
\coordinate (T) at (1.55,0);
\coordinate (Q) at (.5,.525);
\draw[accent!80!black,line width=1.45pt] (M)--(T)--(Q)--cycle;
\fill[accent] (M) circle (1.6pt) node[right,font=\scriptsize] {$M$};
\fill[accent] (T) circle (1.6pt) node[below,font=\scriptsize] {$T$};
\fill[accent] (Q) circle (1.6pt) node[above left,font=\scriptsize] {$Q$};
}
\end{center}
Подставляем три точки в $Ax+By+Cz+D=0$:
\[
\begin{cases}
2A+C+D=0,\\
A+D=0,\\
B+D=0.
\end{cases}
\]
Выберем $D=1$. Тогда $A=-1$, $B=-1$, $C=1$. Получаем эквивалентные записи:
\[
-x-y+z+1=0\quad\Longleftrightarrow\quad\boxed{\beta:\ x+y-z-1=0}.
\]
Нормальный вектор: $\boxed{\vec n_\beta=(1;1;-1)}$.
\subheading{Точная проверка}
\[
M:\ 2+0-1-1=0,\qquad T:\ 1-1=0,\qquad Q:\ 1-1=0.
\]

\heading{5. Скалярное произведение в пространстве}
Для $\vec u=(u_1;u_2;u_3)$ и $\vec v=(v_1;v_2;v_3)$:
\[
\vec u\cdot\vec v=u_1v_1+u_2v_2+u_3v_3,\qquad
|\vec u|=\sqrt{u_1^2+u_2^2+u_3^2}.
\]
Скалярное произведение помогает находить угол, потому что
\[
\vec u\cdot\vec v=|\vec u|\,|\vec v|\cos\theta.
\]

\begingroup\small\setstretch{1.10}
\heading{6. Угол между двумя плоскостями}
\key{Главная идея.} Угол между плоскостями находят через их нормальные векторы. Острый угол $\varphi$ (включая $0^\circ$ и $90^\circ$) вычисляют по модулю скалярного произведения:
\[
\boxed{\cos\varphi=\frac{|\vec n_\alpha\cdot\vec n_\beta|}{|\vec n_\alpha|\,|\vec n_\beta|}},
\qquad 0^\circ\le\varphi\le90^\circ.
\]
\subheading{Пример 3. Продолжение примеров 1 и 2}
Для плоскостей
\[
\alpha:\ x+y-2z=0,\qquad \beta:\ x+y-z-1=0
\]
выбраны нормали $\vec n_\alpha=(1;1;-2)$ и $\vec n_\beta=(1;1;-1)$.

\begin{enumerate}
\item Находим скалярное произведение:
\[
\vec n_\alpha\cdot\vec n_\beta=1\cdot1+1\cdot1+(-2)\cdot(-1)=4.
\]
\item Находим длины нормалей:
\[
|\vec n_\alpha|=\sqrt{1+1+4}=\sqrt6,\qquad
|\vec n_\beta|=\sqrt{1+1+1}=\sqrt3.
\]
\item Подставляем значения в формулу:
\[
\cos\varphi=\frac{|4|}{\sqrt6\,\sqrt3}=\frac4{3\sqrt2}=\boxed{\frac{2\sqrt2}{3}}.
\]
\end{enumerate}
Ответ можно оставить в виде косинуса. При необходимости угол:
\[
\boxed{\varphi=\arccos\frac{2\sqrt2}{3}\approx19{,}47^\circ}.
\]
\heading{7. Что важно проверить перед ответом}
\begin{itemize}[label=$\square$]
\item \textbf{Знаки коэффициентов.} Подставьте все три точки в найденное уравнение.
\item \textbf{Свободный член.} Если плоскость проходит через начало координат, $D=0$.
\item \textbf{Нормаль.} Берите коэффициенты именно при $x,y,z$, без $D$.
\item \textbf{Скалярное произведение.} Умножайте соответствующие координаты и складывайте.
\item \textbf{Длины векторов.} Суммируйте квадраты всех трёх координат.
\item \textbf{Выбор угла.} Для меньшего угла между плоскостями используйте модуль числителя.
\end{itemize}

\endgroup
\clearpage
\thispagestyle{fancy}
\begin{center}{\LARGE\color{darkaccent}Алгоритм: плоскость по трём точкам}\end{center}
\vspace{1mm}
\tikzset{
 fs/.style={draw=accent,rounded corners=2pt,align=center,minimum width=85mm,minimum height=10mm,inner sep=3pt,font=\small},
 decision/.style={diamond,draw=accent,aspect=2.4,align=center,text width=57mm,inner sep=1pt,font=\small},
 terminal/.style={ellipse,draw=darkaccent,align=center,minimum width=52mm,minimum height=9mm,inner sep=2pt,font=\small},
 arr/.style={-{Latex[length=2.3mm]},line width=.75pt,color=darkaccent},
}
\begin{center}
\begin{tikzpicture}[node distance=5mm and 11mm]
\node[terminal] (start) {Даны три точки $P$, $Q$, $R$};
\node[decision,below=of start] (col) {Точки лежат\\на одной прямой?};
\node[fs,right=16mm of col,text width=33mm,minimum width=35mm] (bad) {Единственная плоскость не определяется};
\node[terminal,below=7mm of bad,minimum width=35mm] (badend) {Завершение};
\node[fs,below=of col] (eq) {Запишите $Ax+By+Cz+D=0$.\\Если плоскость проходит через $O$, то $D=0$};
\node[fs,below=of eq] (subst) {Подставьте координаты $P$, $Q$, $R$.\\Решите систему уравнений};
\node[fs,below=of subst] (norm) {Выберите ненулевой общий масштаб\\коэффициентов $A,B,C,D$};
\node[fs,below=of norm] (verify) {Проверьте подстановкой все три точки.\\Нормаль: $\vec n=(A;B;C)$};
\node[decision,below=of verify] (angle) {Требуется угол\\двух плоскостей?};
\node[fs,below=of angle] (twopl) {Найдите нормаль второй плоскости\\и вычислите $\displaystyle\cos\varphi=\frac{|\vec n_1\cdot\vec n_2|}{|\vec n_1|\,|\vec n_2|}$};
\node[terminal,below=of twopl] (end) {Запишите ответ};
\draw[arr] (start)--(col);
\draw[arr] (col)--node[left,font=\scriptsize,fill=white] {нет} (eq);
\draw[arr] (col)--node[above,font=\scriptsize,fill=white] {да} (bad);
\draw[arr] (bad)--(badend);
\draw[arr] (eq)--(subst);
\draw[arr] (subst)--(norm);
\draw[arr] (norm)--(verify);
\draw[arr] (verify)--(angle);
\draw[arr] (angle)--node[left,font=\scriptsize,fill=white] {да} (twopl);
\draw[arr] (twopl)--(end);
\draw[arr] (angle.west)--++(-2.0,0) |- node[pos=.18,left,font=\scriptsize,fill=white] {нет} (end.west);
\end{tikzpicture}
\end{center}
\clearpage

\heading{Тренировочный блок}
\textit{Путь А. 10 заданий от базового уровня к повышенному.} Все координаты указаны в прямоугольной системе координат. Там, где требуется уравнение плоскости, приведите его к виду с целыми коэффициентами без общего делителя.

\subheading{Базовый уровень}
\begin{enumerate}
\item В кубе $ABCDA_1B_1C_1D_1$ с ребром $6$ выбраны оси $Ax$, $Ay$, $Az$ вдоль $AB$, $AD$, $AA_1$. Найдите координаты середин $M$ ребра $BB_1$, $T$ ребра $AB$, $Q$ ребра $AD$.
\item Для точек $P(1;-2;3)$ и $Q(4;0;1)$ найдите координаты вектора $\overrightarrow{PQ}$ и его длину.
\item Даны $\vec u=(2;-1;2)$ и $\vec v=(1;2;0)$. Найдите $\vec u\cdot\vec v$, $|\vec u|$, $|\vec v|$ и угол между векторами.
\end{enumerate}

\subheading{Средний уровень}
\begin{enumerate}[start=4]
\item Для плоскости $3x-y+2z-5=0$ укажите нормальный вектор. Проверьте, принадлежат ли ей точки $X(1;0;1)$ и $Y(1;1;1)$.
\item Составьте уравнение плоскости, проходящей через точки $O(0;0;0)$, $P(2;0;1)$, $Q(0;3;1)$.
\item Составьте уравнение плоскости через точки $P(2;0;0)$, $Q(0;2;0)$, $R(0;0;4)$.
\item Составьте уравнение плоскости через точки $P(3;0;1)$, $Q(1;0;0)$, $R(0;2;0)$. Проверьте результат подстановкой.
\end{enumerate}

\subheading{Повышенный уровень}
\begin{enumerate}[start=8]
\item Найдите косинус меньшего угла между плоскостями $2x-y+2z+1=0$ и $x+2y-2z-3=0$.
\item Найдите меньший угол между плоскостями $x+y-2z=0$ и $x+y+2z-2=0$. Ответ запишите через арккосинус.
\item Плоскость $\alpha$ проходит через $M(2;0;1)$, $T(1;0;0)$, $Q(0;1;0)$. Плоскость $\beta$ проходит через $O(0;0;0)$, $E(2;0;2)$, $F(0;2;1)$. Составьте оба уравнения и найдите косинус меньшего угла между ними.
\end{enumerate}
\clearpage

\heading{Ответы к тренировочному блоку}
\renewcommand{\arraystretch}{1.32}
\begin{tabularx}{\linewidth}{@{}p{11mm}X@{}}
\toprule
\textbf{№}&\textbf{Ответ}\tabularnewline
\midrule
1&$M(6;0;3)$, $T(3;0;0)$, $Q(0;3;0)$.\tabularnewline
2&$\overrightarrow{PQ}=(3;2;-2)$, $|\overrightarrow{PQ}|=\sqrt{17}$.\tabularnewline
3&$\vec u\cdot\vec v=0$, $|\vec u|=3$, $|\vec v|=\sqrt5$, угол $90^\circ$.\tabularnewline
4&$\vec n=(3;-1;2)$; $X$ принадлежит, $Y$ не принадлежит плоскости.\tabularnewline
5&$3x+2y-6z=0$.\tabularnewline
6&$2x+2y+z-4=0$.\tabularnewline
7&$2x+y-4z-2=0$.\tabularnewline
8&$\displaystyle\cos\varphi=\frac49$.\tabularnewline
9&$\displaystyle\varphi=\arccos\frac13\approx70{,}53^\circ$.\tabularnewline
10&$\alpha:\ x+y-z-1=0$; $\beta:\ 2x+y-2z=0$; $\displaystyle\cos\varphi=\frac{5}{3\sqrt3}=\frac{5\sqrt3}{9}$.\tabularnewline
\bottomrule
\end{tabularx}

\heading{Итоговая памятка}
\begin{tabularx}{\linewidth}{@{}p{47mm}X@{}}
\toprule
\textbf{Сигнал в задаче}&\textbf{Что делать}\tabularnewline
\midrule
Даны вершины куба&Выбрать начало координат и оси вдоль рёбер.\tabularnewline
Дана середина ребра&Найти средние арифметические координат концов.\tabularnewline
Даны три точки плоскости&Подставить их координаты в $Ax+By+Cz+D=0$.\tabularnewline
Нужно найти нормаль&Взять тройку коэффициентов $(A;B;C)$.\tabularnewline
Нужен угол плоскостей&Найти две нормали и применить формулу для $\cos\varphi$.\tabularnewline
Получено уравнение&Проверить подстановкой все заданные точки.\tabularnewline
\bottomrule
\end{tabularx}
\vspace{8mm}
\begin{center}
{\small\color{darkaccent}Лёвин Артём Александрович}\par
{\small\color{muted}эксклюзивно для Матвея Горбачёва}
\end{center}
\end{document}
